\ifdefined\isincluded
\else
\documentclass[a4paper, 12pt]{report}
\usepackage[margin=0.8in]{geometry}
\usepackage{amsmath}
\usepackage{amsfonts}
\usepackage{amssymb}
\usepackage{graphicx}
\usepackage[english]{babel}
\usepackage[autostyle, english = american]{csquotes}
\MakeOuterQuote{"}
\usepackage{setspace}
\onehalfspacing
\usepackage{booktabs}
\usepackage{longtable}
\usepackage{tabularx}
\usepackage{array}
\usepackage{url}
\usepackage{cite}

\makeatletter
\let\oldincludegraphics\includegraphics
\renewcommand{\includegraphics}[2][]{%
  \IfFileExists{#2}{%
    \oldincludegraphics[#1]{#2}%
  }{%
    \IfFileExists{images/#2}{%
      \oldincludegraphics[#1]{images/#2}%
    }{%
      \fbox{\parbox[c][3cm][c]{0.8\linewidth}{\centering \texttt{[Image: #2]}}}%
    }%
  }%
}
\makeatother

\begin{document}
\chapter{Results and Discussions}
\fi

\label{chap.5}
\section{Quantitative Benchmark Performance Evaluation}
This chapter uses empirical data and in-depth discussions from 912 benchmark evaluations on the stateful $\tau^2$-Bench Retail domain to compare Multi-Round Text ReAct to Recursive Latent ReAct across 114 jobs assessed at recursion levels $r \in \{0, 1, 2, 3\}$.

Table~\ref{tab:macro_performance} summarizes the benchmark performance for each of the 912 execution runs. We distinguish clearly between: (i) the continuous goal satisfaction accuracy ($S_{\text{goal}} \times 100\%$) averaged over all recursion depths ($r \in \{0, 1, 2, 3\}$)---reported as an observable descriptive difference of 50.9\% for Latent ReAct compared to 50.3\% for Text ReAct (+0.6 pp), increasing to 56.4\% as opposed to 54.4\% at terminal depth ($r=3$, a gain of 2.0 pp); (ii) Out of 114 tasks, 64 tasks (56.1\%) for Latent ReAct and 62 tasks (54.4\%) for Text ReAct achieved the general binary task success rate ($S_{\text{goal}} \ge 0.70$) at terminal depth; and (iii) 33 tasks for Latent ReAct and 30 tasks for Text ReAct achieved clean terminal success ($S_{\text{goal}} \ge 0.70$ without tool exceptions or policy violations). Latent ReAct performs better than Text ReAct at the task level, with an error recovery yield of 33 initially failing tasks. (28.9\% across all 114 benchmark tasks) over the course of future recursive passes, compared to 28 tasks (24.6\%). Latent ReAct improves goal accuracy on hard processes by 4.9 percentage points (50.8\% versus 45.9\%) and decreases enterprise policy violations by 23.1\% (10 versus 13 jobs). Latent ReAct simultaneously achieves a 25.2\% decrease in generated tokens, a 15.6\% reduction in total inference computation, and an observed 26.1\% reduction in mean end-to-end latency (8.35 seconds against 11.30 seconds) on the examined benchmark environment.

\begin{table}[h!]
\centering
\small
\caption{Overall Benchmark Performance Summary for 912 Execution Runs on $\tau^2$-Bench Retail Domain.}
\label{tab:macro_performance}
\vspace{4pt}
\resizebox{\linewidth}{!}{%
\begin{tabular}{lccc}
\toprule
\textbf{Performance Dimension / Metric} & \textbf{Text ReAct} & \textbf{Latent ReAct} & \textbf{Relative Advantage} \\
\midrule
Mean Goal Accuracy ($r=0\dots 3$) & 50.3\% & \textbf{50.9\%} & \textbf{+0.6 pp (Latent)} \\
Terminal Goal Accuracy ($r=3$) & 54.4\% & \textbf{56.4\%} & \textbf{+2.0 pp (Latent)} \\
Terminal Binary Success ($S_{\text{goal}} \ge 0.70$) & 54.4\% (62 tasks) & \textbf{56.1\% (64 tasks)} & \textbf{+1.7 pp (Latent)} \\
Hard Workflow Goal Accuracy & 45.9\% & \textbf{50.8\%} & \textbf{+4.9 pp (Latent)} \\
Clean Success Tasks (Terminal $r=3$, out of 114) & 30 tasks (26.3\%) & \textbf{33 tasks (28.9\%)} & \textbf{+10.0\% (Latent)} \\
Benchmark Error Recovery Yield (out of 114) & 24.6\% (28 tasks) & \textbf{28.9\% (33 tasks)} & \textbf{+17.9\% (Latent)} \\
Policy Constraint Violations & 13 tasks & \textbf{10 tasks} & \textbf{23.1\% Reduction (Latent)} \\
Mean End-to-End Latency & 11.30 s & \textbf{8.35 s} & \textbf{26.1\% Observed Reduction (Latent)} \\
Total Output Generated Tokens & 166,779 tokens & \textbf{124,805 tokens} & \textbf{25.2\% Savings (Latent)} \\
Total Inference Computation & 13,076 TFLOPs & \textbf{11,038 TFLOPs} & \textbf{15.6\% Savings (Latent)} \\
Isolated Tool Selection Accuracy & \textbf{57.5\%} & 54.2\% & -3.3 pp (Latent - Text) \\
\bottomrule
\end{tabular}%
}
\end{table}

Table~\ref{tab:depth_progression} displays the sequential evolution of task goal satisfaction accuracy, latency, tool precision, and token consumption spanning recursion levels zero to three. Both architectures show how recursive deliberation works by raising target accuracy from about 45\% at round zero to over 54\% at round three. On the other hand, Latent ReAct scales with a far higher computational efficiency, finishing round three in 13.83 seconds as opposed to Text ReAct's 18.56 seconds.
\begin{table}[h!]
\centering
\small
\caption{Benchmark Evaluation Breakdown by Recursion Depth ($r = 0, 1, 2, 3$), with a mean continuous accuracy of goal satisfaction ($S_{\text{goal}} \times 100\%$).}
\label{tab:depth_progression}
\vspace{4pt}
\resizebox{\linewidth}{!}{%
\begin{tabular}{lcccccccc}
\toprule
& \multicolumn{4}{c}{\textbf{Text ReAct (Baseline)}} & \multicolumn{4}{c}{\textbf{Recursive Latent ReAct (Ours)}} \\
\cmidrule(lr){2-5} \cmidrule(lr){6-9}
\textbf{Depth} & \textbf{Goal Acc (\%)} & \textbf{Lat (s)} & \textbf{Tool (\%)} & \textbf{Tokens} & \textbf{Goal Acc (\%)} & \textbf{Lat (s)} & \textbf{Tool (\%)} & \textbf{Tokens} \\
\midrule
$r = 0$ & 45.7 & 4.36 & 58.8 & 16,891 & 45.3 & \textbf{3.05} & 56.1 & \textbf{11,719} \\
$r = 1$ & 48.2 & 8.82 & 57.9 & 33,442 & \textbf{48.5} & \textbf{6.52} & 52.2 & \textbf{25,077} \\
$r = 2$ & 52.8 & 13.45 & 57.3 & 49,412 & \textbf{53.3} & \textbf{9.98} & 55.0 & \textbf{37,387} \\
$r = 3$ & 54.4 & 18.56 & 56.1 & 67,034 & \textbf{56.4} & \textbf{13.83} & 53.5 & \textbf{50,622} \\
\bottomrule
\end{tabular}%
}
\end{table}

Table~\ref{tab:complexity_tier} looks at the terminal performance and latency scaling achieved at full recursion depth ($r=3$) across job complexity tiers. The cohort's single-run execution latencies at $r=3$ include the total inference time over four full deliberation runs, with an average of 18.56 s for Text ReAct and 13.83 s for Latent ReAct (see Table~\ref{tab:depth_progression}). On the other hand, the grand mean latencies shown in Table~\ref{tab:macro_performance} (11.30 s for Text ReAct vs. 8.35 s for Latent ReAct) indicate the average operating cost for all 912 runs over all recursion depths ($r \in \{0, 1, 2, 3\}$).10.18 s vs. 7.71 s on Easy, 11.71 s vs. 8.36 s on Medium, and 12.00 s vs. 8.97 s on Hard are the tier-specific averages for all depths. On hard processes involving multi-entity coordination, replacement checks, and restocking policy deductions, Latent ReAct performs noticeably better than Text ReAct (50.8\% vs. 45).Text ReAct maintains a little advantage over Medium single-mutation processes, with 9\% accuracy and 24.5\% faster execution at $r=3$.

\begin{table}[h!]
\centering
\small
\caption{Task Complexity Tier Performance and Execution delay at Terminal Depth $r=3$ (single-run execution delay and reported continuous goal fulfillment accuracy $S_{\text{goal}} \times 100\%$).}
\label{tab:complexity_tier}
\vspace{4pt}
\resizebox{\linewidth}{!}{%
\begin{tabular}{lcccc}
\toprule
& \multicolumn{2}{c}{\textbf{Text ReAct}} & \multicolumn{2}{c}{\textbf{Recursive Latent ReAct}} \\
\cmidrule(lr){2-3} \cmidrule(lr){4-5}
\textbf{Complexity Tier} & \textbf{Terminal Goal Acc (\%)} & \textbf{Terminal Latency (s)} & \textbf{Terminal Goal Acc (\%)} & \textbf{Terminal Latency (s)} \\
\midrule
Easy Tier (38 Tasks) & 53.8\% & 16.67 s & \textbf{57.5\%} & \textbf{12.89 s} \\
Medium Tier (38 Tasks) & \textbf{63.7\%} & 19.54 s & 60.9\% & \textbf{13.93 s} \\
Hard Tier (38 Tasks) & 45.9\% & 19.46 s & \textbf{50.8\%} & \textbf{14.68 s} \\
\bottomrule
\end{tabular}%
}
\end{table}

\section{Behavioral Accuracy and Computational Scaling}
Figure~\ref{fig:acc_depth} displays the evolution of overall task completion accuracy across recursion rounds $r \in \{0, 1, 2, 3\}$. Both architectures exhibit monotonic performance gains as recursion depth rises from $r=0$ to $r=3$. In contrast, Latent ReAct exhibits a faster accuracy trend from 45.3 percent at round zero to 56.4 percent at round three, outperforming Text ReAct by an observed two percentage points at full recursion depth. In complex workflows, the accumulation of verbose textual self-critique within the prompt context may lead to model fixation on early errors, potentially creating needless token history. The idea that Latent ReAct maintains relational problem representations in continuous space, avoiding cumulative textual history bloat and permitting multi-step error recovery, is supported by our recovery yield measurements.
Figure~\ref{fig:flops_depth} displays the scaling of mean computational inference floating-point operations per task run throughout recursive depths $r \in \{0, 1, 2, 3\}$. It is important to distinguish between the per-run operational scale plotted in Figure~\ref{fig:flops_depth} and the benchmark-wide totals shown in Table~\ref{tab:macro_performance}. The total cumulative sum of all floating-point operations throughout all 456 execution runs on the 114 benchmark tasks is the stated computational volume (13,076 TFLOPs for Text ReAct versus 11,038 TFLOPs for Latent ReAct) in Table~\ref{tab:macro_performance}. The presented curves in Figure~\ref{fig:flops_depth} illustrate the average computational cost per individual task execution at each discrete recursion depth. The per-run compute scales from seven point eight nine TFLOPs for Text ReAct and seven point zero seven TFLOPs for Latent ReAct at round zero ($r=0$) to thirty-four point nine four versus twenty-nine point three six TFLOPs at round two ($r=2$) and fifty-two point eight four TFLOPs for Text ReAct versus forty-three point nine two TFLOPs for Latent ReAct at round three ($r=3$). This results in a sixteen percent compute reduction per task at full recursion depth. The autoregressive creation of explicit natural language tokens required by Text ReAct's intermediate stages results in a significant increase in computation per run. Latent ReAct uses matrix multiplications in lightweight neural adapters for all intermediate deliberation stages, avoiding token generation overhead and achieving a 15.6\% reduction in cumulative benchmark computation.

\begin{figure}[htbp]
\centering
\begin{minipage}[t]{0.48\textwidth}
\centering
\includegraphics[width=\linewidth]{accuracy_vs_recursive_depth.png}
\caption{Total task completion accuracy as a function of recursive depth ($r = 0, 1, 2, 3$) for 114 corporate jobs comparing Multi-Round Text ReAct to Recursive Latent ReAct.}
\label{fig:acc_depth}
\end{minipage}\hfill
\begin{minipage}[t]{0.48\textwidth}
\centering
\includegraphics[width=\linewidth]{flops_vs_recursive_depth.png}
\caption{Table~\ref{tab:macro_performance} shows the average computational inference expenditure per task run (TFLOPs) over recursion depths ($r \in \{0, 1, 2, 3\}$), comparing single-run scaling curves with benchmark-wide cumulative totals.}
\label{fig:flops_depth}
\end{minipage}
\end{figure}

\section{Geometric State Convergence and Dimensionality Results}
Figure~\ref{fig:geom_convergence} displays the empirical step-wise state drift across recursion passes along with 95\% bootstrap confidence intervals. The visual trajectory experimentally validates finding 1 from our geometric analysis: Latent ReAct concentrates 64\% of its total displacement in the initial transition ($r=0 \to 1$), falling sharply from $0.065$ to $0.018$ before stabilizing. Text ReAct, on the other hand, exhibits a drift that is approximately constant, averaging $0.034$ every step ($0.038$, $0.032$, and $0.033$ across consecutive passes). A paired Wilcoxon signed-rank test over all 114 tasks shows that this stepwise decline from first to subsequent transitions is highly significant ($p = 1.9 \times 10^{-20}$). Even after removing tasks with duplicate states, the latent contraction remains strong ($p = 7.1 \times 10^{-15}$), whereas the text drift is flat ($p = 0.96$). Technically, we do not provide a guaranty of fixed-point contractivity with the unconstrained residual adapters, but we do observe empirically a substantial stabilization of the trajectory toward a compact manifold, similar to convergence. We find that trajectory stabilization is an empirical geometric pattern, rather than a direct, verified cause of task completion.

Figure~\ref{fig:geom_dimensionality}: Participation Ratio measured over two hundred bootstrap subsets for the effective dimensionality of the representation manifold. In each round latent ReAct has a smaller, lower dimensional subspace than Text ReAct ($18.9 \to 15.2$ dimensions for latent vs. $22.0 \to 16.8$ dimensions for text). Both architectures achieve 20--23\% compression over recursive cycles. This is consistent with the idea that latent communication, by removing unnecessary linguistic variability, preserves a compact, non-collapsing manifold that is suitable for enterprise thinking.

\begin{figure}[htbp]
\centering
\begin{minipage}[t]{0.48\textwidth}
\centering
\includegraphics[width=\linewidth]{fig_A_convergence.png}
\caption{95\% bootstrap confidence intervals for stepwise representation drift over recursive passes.}
\label{fig:geom_convergence}
\end{minipage}\hfill
\begin{minipage}[t]{0.48\textwidth}
\centering
\includegraphics[width=\linewidth]{fig_B_dimensionality.png}
\caption{Over 200 bootstrap subsets were used to evaluate the participation ratio and effective dimensionality across recursive passes.}
\label{fig:geom_dimensionality}
\end{minipage}
\end{figure}

The effect sizes (Cliff's $\delta$ with 95\% bootstrap confidence intervals) for each of the five possible trajectory metrics that compare clean successes to failures are shown in Figure~\ref{fig:geom_outcome}. Across the family of statistical hypotheses, all intervals embrace zero after Holm-Bonferroni correction (minimum adjusted $p = 0.54$), empirically proving that downstream transaction success is not only determined by macroscopic representation geometry.

\begin{figure}[htbp]
\centering
\includegraphics[width=0.48\linewidth]{fig_C_outcome_link.png}
\caption{A null connection with global representation geometry was confirmed by comparing clean successes against failures using effect sizes (Cliff's $\delta$ with 95\% bootstrap confidence intervals) across five possible trajectory features.}
\label{fig:geom_outcome}
\end{figure}

\section{Discussion of the Tool Selection Trade-Off}
In single-step isolated tool selection, Text ReAct's observed descriptive accuracy was 57.5\%, whereas Latent ReAct's was 54.2\%. This shows a 3.3 percentage point negative difference (Latent minus Text; bootstrap 95 percent confidence interval of $[-5.1\%, -1.5\%]$). This observed variance reflects the superficial lexical advantage of Text ReAct, where clear API names and schema definitions are instantly accessible in the prompt strings. Text ReAct's lower ultimate accuracy of 45 could be explained in one way.9\% on Hard tasks is that unrestrained verbal reflection finds it difficult to correct compounding relational faults on complicated workflows when there is no ongoing trajectory stabilization. On the other hand, continual latent feedback seems to preserve relational consistency during deliberation passes, enabling Latent ReAct to achieve higher terminal accuracy and resolve 10\% more corporate jobs cleanly.

\ifdefined\isincluded
\else
\end{document}
\fi
